\textit{Función de fitness (OneMax):} $f(x)=\sum_{i=1}^{6} x_i$

\textit{Población inicial}

\begin{center}
\begin{tabular}{c|c}
Individuo & Fitness \\
\hline
000000 & 0 \\
111010 & 4 \\
010100 & 2 \\
000001 & 1 \\
\end{tabular}
\end{center}

Se toman los 3 mejores:

\[
\{111010(4),\;010100(2),\;000001(1)\}
\]


\bigskip
\textbf{Iteración 1}

- La recombinación usa mitad del primero + mitad del segundo (3 bits y 3 bits).

\[
(1,2):\;111010 + 010100
\]

\[
111|100 = 111100
\]

\[
(2,3):\;010100 + 000001
\]

\[
010|001 = 010001
\]

Nueva población:

\begin{center}
\begin{tabular}{c|c}
Individuo & Fitness \\
\hline
111100 & 4 \\
010001 & 2 \\
111010 & 4 \\
010100 & 2 \\
000001 & 1 \\
000000 & 0 \\
\end{tabular}
\end{center}

Mejores 3:

\[
\{111100(4),\;111010(4),\;010100(2)\}
\]

\bigskip
\textbf{Iteración 2}

\[
(1,2):\;111100 + 111010
\]

\[
111|010 = 111010
\]

\[
(2,3):\;111010 + 010100
\]

\[
111|100 = 111100
\]

Nueva población:

\begin{center}
\begin{tabular}{c|c}
Individuo & Fitness \\
\hline
111010 & 4 \\
111100 & 4 \\
111100 & 4 \\
111010 & 4 \\
010100 & 2 \\
010001 & 2 \\
000001 & 1 \\
000000 & 0 \\
\end{tabular}
\end{center}

Mejores 3:

\[
\{111100(4),\;111010(4),\;111100(4)\}
\]

\bigskip
\textbf{Iteración 3}

\[
(1,2):\;111100 + 111010
\]

\[
111|010 = 111010
\]

\[
(2,3):\;111010 + 111100
\]

\[
111|100 = 111100
\]

Nueva población:

\begin{center}
\begin{tabular}{c|c}
Individuo & Fitness \\
\hline
111010 & 4 \\
111100 & 4 \\
111100 & 4 \\
111010 & 4 \\
111100 & 4 \\
111010 & 4 \\
010100 & 2 \\
010001 & 2 \\
000001 & 1 \\
000000 & 0 \\
\end{tabular}
\end{center}

\bigskip
En la población final existen varios individuos que alcanzan el valor máximo de fitness observado. 
Por lo tanto, cualquier elemento de $\arg\max f(x)$ dentro de la población constituye una solución válida. 
En este caso, el mayor valor alcanzado es $f(x)=4$.

Sin pérdida de generalidad, podemos seleccionar uno de estos individuos como candidato final.

\textbf{Candidato final}

\[
\boxed{111100}
\]

\[
f(111100)=4
\]
